\section{Analysis}
\label{sec:analysis}

\subsection{Most submissions change how the run goes, not how the model learns}
\label{sec:families}

Everything so far has been a score, and a score says only whether a submission worked. It does not
say whether the agent designed an algorithm or tuned one that was already there, and those are the
two outcomes this paper exists to separate. Counting lines does not separate them either: a one-line
diff can replace a learning rule, and a thousand-line refactor can leave the training procedure
exactly as it was. The distinction has to be read off the submitted code, by asking which part of the
training procedure it modifies.

% Generated by scripts/make_families_table.py -- do not edit by hand.
\begin{table}[t]
\centering
\caption{\textbf{What the submissions change.} The 263 submissions that changed anything that could be classified, out of 280 (\S\ref{sec:setup}); Kimi~K3 is outside this corpus. Families are not exclusive --- a submission matches $3.13$ of them on average --- so the shares do not sum to one, and the line that matters is how many submissions reach the learning side at all. The lower panel gives that share by reasoning effort.}
\label{tab:families}
\small
\begin{adjustbox}{max width=\linewidth}
\begin{tabular}{llrr}
\toprule
 & \textbf{Family} & \textbf{$n$} & \textbf{share} \\
\midrule
\multirow{4}{*}{\rotatebox{90}{\footnotesize run}}
& how long it trains, how often it saves & 253 & 96.2\% \\
& the training hyperparameters & 195 & 74.1\% \\
& which checkpoint to keep & 105 & 39.9\% \\
& how much trainable capacity, and where & 73 & 27.8\% \\
\midrule
\multirow{4}{*}{\rotatebox{90}{\footnotesize learning}}
& the loss it optimizes & 87 & 33.1\% \\
& the supervision it learns from & 66 & 25.1\% \\
& the update rule itself & 23 & 8.7\% \\
& the data it trains on & 21 & 8.0\% \\
\midrule
& \textbf{any learning family} & \textbf{122} & \textbf{46.4\%} \\
& \emph{run side only} & 141 & 53.6\% \\
\bottomrule
\end{tabular}
\end{adjustbox}
\par\vspace{5pt}
\begin{adjustbox}{max width=\linewidth}
\begin{tabular}{lrrrrrr}
\toprule
\emph{reaching the learning side} & \texttt{none} & \texttt{low} & \texttt{medium} & \texttt{high} & \texttt{xhigh} & \texttt{max} \\
\midrule
share of submissions & 8\% & 39\% & 33\% & 49\% & 65\% & 64\% \\
\bottomrule
\end{tabular}
\end{adjustbox}
\end{table}


We group every change into eight families on two sides of that line (Table~\ref{tab:families}). The
grouping is assigned by a separate language model reading each submitted diff against the definitions
that follow. Four
change \emph{how this run goes}: how long it trains and how often it saves; the training
hyperparameters, such as the learning rate or the batch size; which of the checkpoints it produced to
keep; and how much trainable capacity to attach and where, such as the rank and
placement of an adapter. Four change \emph{how the model learns}: the loss it optimizes, by adding,
removing or reweighting a term; the supervision it learns from, by introducing a signal the procedure
did not have before; the update rule itself, replaced by a different one; and the data the procedure
trains on. The families are not exclusive --- a submission matches $3.13$ of them on
average, since changing a loss usually drags a hyperparameter along with it --- so we report how many
submissions reach a side rather than assigning each to a single family.

Of the 280 submissions, 17 made no change that could be classified. Of the remaining 263, 141 stay
entirely on the run side and only 122 touch how the model learns. Four hours, a whole repository, and
a task statement that says in as many words to improve this training algorithm, and more than half of
the submissions never reach that layer.

% All figures here are means of the mapped score of \S\ref{sec:scoring}, computed
% from build/lark/mapped_table.csv joined to build/family_rows.jsonl. Reported as
% means rather than as win rates: over ten tasks a win rate is too unstable to
% carry a comparison, and it discards the margin, which is the quantity at issue.
Do the submissions that reached it do better? On the scale of \S\ref{sec:scoring} they do, and by a
wide margin: submissions that touch the learning procedure average $0.226$ against $0.126$ for those
that stay on the run side, a gap of $0.100$ against a standard error of $0.022$. It is not the
artifact of a single task --- dropping agentic RL, where imitation learning lifts the whole column,
still leaves $0.182$ against $0.128$ --- nor of a single system, since the ordering holds within four
of the five models in this corpus. It is also not a randomized comparison: the systems that reach the
learning procedure more often are the stronger ones to begin with, so the gap is the difference
between the submissions that go there and the submissions that do not, rather than the effect of
going there.

Read together, the two numbers say that the algorithmic layer is where the distance actually gets
closed, and that most submissions never go to it. What separates a submission that reaches that layer
from one that does not is not effort spent but a step taken first: reading the training dynamics as a
specific failure mechanism, and then addressing that mechanism.

\subsection{Reasoning effort buys nerve, and nerve is what pays}
\label{sec:effort}

Is there anything that pushes agents down to that layer? The setup has exactly one knob that can be
turned on its own, the reasoning-effort level, and taking it from the lowest setting to the highest
shows clearly what it buys and what it does not.

It buys nerve. The share of submissions that touch the learning algorithm rises from $8.0\%$ to
$64.0\%$. At low effort the submissions move
budgets, logging and optimization knobs; at high effort they operate on objectives, replace learning
rules, and add supervision to the procedure.

It buys attempts. Within the Codex grid, the only one that exposes the lowest setting, the median
configuration goes from 4 evaluations inside its four hours to 16, from
18 edited lines to 246, and from $11$k output tokens to $109$k. Cost follows: the median exploration
cost per task rises from \$1.69 to \$34.60. The exploration stage of the whole
evaluation consumed \$5{,}334 of API calls, with Kimi~K3 converted from Chinese yuan and with neither
the GPU hours of the twelve-hour runs nor the evaluator's compute included.

It buys completion. Of the 19 cells that score zero, 8 ended their four hours without a usable patch
--- four of them workspaces the host found empty after the agent exited early --- and 11 submitted a
complete patch that started its twelve-hour run but finished with nothing satisfying the contract,
most often no loadable merged model written to disk. All 19 terminated normally: the failure is in what
the agent submitted. They concentrate at low effort, with 12 of the 19 at the two lowest levels and
one each at the two highest.

% Means of the mapped score from build/lark/mapped_table.csv. The Codex figures
% hold the harness fixed, which is why they are quoted beside the pooled ones.
And it does buy a result, though a small one in absolute terms. On the scale of \S\ref{sec:scoring}
the mean rises from $0.094$ at the lowest level to $0.196$ at the highest, a gap of $0.102$ against a
standard error of $0.027$; within the Codex grid, where the harness is held fixed, it rises at every
step, $0.094$ to $0.204$. Set beside the eightfold rise in submissions that reach the learning
procedure, the score roughly doubles, and $0.196$ is still only a tenth of the way from the shipped
algorithm to the optimum. Reasoning effort works, then, by making an agent attempt the thing that
pays rather than by making the attempt itself better: it brings more agents to the loop that
matters --- reading the training dynamics, naming the mechanism that is failing, changing it --- and
what they gain is about what arriving there is worth.

\subsection{What the submissions that reached the algorithmic layer did}
\label{sec:reached}

Among the 122 submissions that touched how the model learns, a few changed not one step of the
procedure but what the task was taken to be. Three are worth reading in full, and they come from
three different tasks.

\textbf{Turning a task that does no training into one that does.} One-shot pruning is defined to do
exactly one thing: score which weights to remove, remove them once, stop. The repository's own
procedure leaves a perplexity of $53.4$. One submission replaced it with a three-stage pipeline ---
a different rule for selecting and updating the surviving weights, then a round of layerwise
distillation, then a masked knowledge-distillation fine-tune of the whole model (AdamW, 666 steps,
cosine decay) --- and brought the perplexity to just over $13$. Its notes record a diagnosis along
the way: a first attempt scored $572$, absurdly bad, because the weight-allocation step propagated
activations forward through the network and overwrote layer 0's input in place, so the pruning step
was reading layer 31's activations.

\textbf{Turning a closed form into an optimization problem.} Weight averaging ships a uniform mean
over 72 candidate models. One submission first built itself an instrument: the relevant tensors of
all 72 models packed into a single matrix resident in GPU memory and the proxy images preprocessed
and held, so that loading a coefficient vector and scoring it became one matrix multiplication and
one forward pass --- $0.38$ seconds, against roughly $190$ before. On that instrument it ranked five
methods, all measured on its own rig: best single model $0.6935$, uniform average $0.6880$, top-$k$
by accuracy $0.6945$, greedy soup $0.7025$, and coefficients learned directly by cross-entropy with
Adam $0.7020$, the last two tied. It also recorded two routes that did not work: extrapolating along
a single direction collapses accuracy, and a logit-ensemble proxy does not rank candidates reliably.

\textbf{Replacing reinforcement learning with imitation learning.} Multi-turn agentic RL ships GRPO. The
submissions that reached a perfect score judged that on this task it pays to learn from the optimal
solution first: generate boards in quantity, label every step with its optimal move, and fine-tune on
that supervision; one went further with \textsc{dagger}, letting the policy walk and adding the
correct answer wherever it went.

The three have one thing in common. Each built something measurable before acting: a solver to
establish the task's ceiling, an evaluation rig five hundred times faster than the one it was given,
a localisation of which layer's activations were being overwritten. This is exactly the capability
\S\ref{sec:families} finds missing --- and among 263 submissions it is the exception.
